\section{研究背景与核心动机}\label{sec:background}

\begin{tcolorbox}[hypothesisstyle,
  title={\normalsize\bfseries\color{natgreen}PrimeCoDec 故事脉络 \quad {\normalfont\itshape Story Narrative}}]
\textbf{一、问题根源}——纠错码在无线通信、光通信、数据存储、量子计算、存算一体等\textbf{所有信息系统}中不可或缺，但当前设计范式仍停留在``专家手工作坊''模式：过度依赖专家经验、优化严重依赖大规模仿真、算法与硬件深度割裂——三大困境导致ECC设计高度低效，严重制约新兴信息系统的快速部署。

\vspace{6pt}
\textbf{二、框架主张}——\textbf{PrimeCoDec} 提出：用户仅输入原始需求，框架即可全自动完成从纠错码编码设计、译码算法选择、硬件实现到行业标准输出的全链路自动化，打破ECC设计中一切专家壁垒。

\vspace{6pt}
\textbf{三、技术支柱}——四个均衡的技术层次共同支撑框架：
\begin{itemize}[nosep]
  \item \textbf{编码层}：AI驱动的码型自动搜索与构造（跨域通用信道表示）
  \item \textbf{译码层}：神经算子替代蒙特卡洛仿真，实现快速算法评估与调优
  \item \textbf{硬件层}：算法-硬件联合代理优化，支持共模架构自动实现
  \item \textbf{标准层}：LLM驱动的需求形式化与行业标准草案自动生成
\end{itemize}

\vspace{6pt}
\textbf{四、核心验证}——以\textbf{5G NR标准（3GPP Release 17）}为基准，验证PrimeCoDec的全链路能力：从自然语言需求输入，到Polar码/LDPC码最优方案、共模译码器硬件，最终输出修订后的3GPP标准草案——一个完整的从需求到标准的自动化闭环。
\end{tcolorbox}

\vspace{8pt}
\subsection{纠错码：横跨全域信息系统的核心基础技术}

纠错码技术诞生于1948年Shannon信息论，历经七十余年发展，已成为数字信息基础设施不可或缺的核心组件——其触角延伸至当今几乎所有信息系统，是横跨无线通信、光通信、数据存储、新型计算乃至量子信息的\textbf{共性底层技术}。

\textbf{无线通信}方面，在5G NR标准（3GPP Release 15+）中，LDPC码\cite{gallager1962}承担eMBB场景数据信道编码任务，Polar码\cite{arikan2009}被采纳为控制信道方案。下一代6G系统对URLLC（超可靠低时延通信）提出了更严苛的编码需求——在更短码长（$n < 256$）与更高可靠性（BLER $\leq 10^{-9}$）约束下保障极端低时延。

\textbf{光通信}方面，软判决前向纠错（SD-FEC）已广泛采用LDPC码，实现净编码增益（NCG）超过10~dB，成为400G/800G乃至未来Tb/s光传输系统的关键使能技术。

\textbf{存储系统}方面，随着NAND Flash存储密度从SLC向QLC跃升，信道特性呈现显著的非高斯性与时变性，亟需针对磨损非平稳噪声的定制化ECC方案。

\textbf{量子计算}方面，量子纠错码（QECC）是实现容错量子计算的唯一已知途径\cite{shor1995}。以表面码\cite{fowler2012}为代表的拓扑量子纠错方案高度依赖对量子噪声模型的精确建模与高效经典译码。

\textbf{存算一体}方面，忆阻器（Memristor）阵列的模拟电导分布噪声构成独特的存储信道模型，现有ECC体系缺乏针对此类新型器件的系统性设计工具。

上述五大应用领域共同指向一个深刻事实：\textbf{只要存在信息的传输与存储，纠错码就无处不在}。然而令人震惊的是，面向如此广泛且多样化的应用，当前ECC的设计方式却几乎停留在香农时代的手工作坊模式。

\subsection{现有ECC设计范式的三大结构性困境}\label{sec:pain_points}

纵观上述所有应用场景，ECC设计者每次面对新需求时，都不得不从头开始应对同样的三大系统性困境：

\textbf{困境一：码型优化与算法设计高度依赖专家先验知识。}LDPC基图设计需要深厚的环分析、度分布优化与密度进化理论基础\cite{richardson2001}；Polar码编码序列的选择依赖Tal--Vardy算法\cite{tal2013}等复杂理论工具；译码算法的参数调优（迭代次数、量化位宽、调度策略）需要多年工程经验积累。这一\textbf{专家壁垒}严重制约了ECC技术向新型应用场景的快速推广，也使得每当出现新系统、新信道时，ECC方案的从头设计成为不可避免的高成本工程。

\textbf{困境二：性能评估与参数优化严重依赖大规模仿真。}获取一条可靠的BER曲线需要在每个目标SNR点累计数亿次随机噪声样本。以码长$n = 10^4$、目标BER为$10^{-8}$的LDPC码为例，完整评估需约$10^{12}$次信道仿真，即便借助高端GPU并行加速也需数小时乃至数天。这一瓶颈使得大规模码型参数搜索几乎不可行，\textbf{大量人力与计算资源被耗费在重复仿真而非真正的创造性设计上}。

\textbf{困境三：算法优化与硬件实现深度脱节。}算法工程师在优化BER性能时通常不考虑译码器的硬件实现代价；当算法确定后转入硬件实现阶段，往往发现在面积/功耗约束下无法直接实现，导致大量推倒重来。更深层的问题是：\textbf{当前没有任何框架能够将算法优化与硬件部署、乃至行业标准制定统一于同一闭环}，每个环节都是孤立的"信息孤岛"。

这三大困境共同造成了ECC领域高度低效的设计生态：仅有极少数行业（如5G NR无线通信）凭借数十年专家深度投入，才建立起完备且充分优化的行业标准（3GPP TS 38.212等\cite{3gpp38212}）；绝大多数新兴应用场景（量子计算、忆阻器存算一体、6G新波形等）至今仍缺乏系统性的ECC设计支撑。

\subsection{人工智能带来的变革机遇}\label{sec:ai_opportunity}

近年来，AI技术在若干关键方向取得的突破为克服上述三大困境、实现ECC设计全链路自动化提供了历史性机遇：

\textbf{图神经网络（GNN）用于结构化学习。}因子图是LDPC码、Polar码等现代纠错码的统一数学表示工具，GNN天然适合对因子图的拓扑结构进行编码学习，已有大量工作证明其在神经译码\cite{nachmani2018,kang2022}领域的有效性。

\textbf{神经算子（Neural Operator）用于泛函逼近。}以DeepONet\cite{lu2021}与FNO\cite{kovachki2023}为代表的神经算子，已在PDE求解、湍流预测等领域展现出强大的泛化能力，为脱离蒙特卡洛仿真提供了全新理论工具。

\textbf{大语言模型（LLM）用于需求理解与代码生成。}GPT-4o、Claude等LLM已展示出强大的自然语言理解与代码生成能力，在EDA领域的初步应用\cite{liu2023,chang2024}表明RTL代码生成与硬件需求理解的可行性。

\textbf{硬件代理模型（Hardware Surrogate Models）。}以GNN与MLP为核心的轻量硬件性能代理模型\cite{ustun2020,geng2021}在逻辑综合结果预测方面展现出接近综合工具精度的同时实现了数量级的速度提升。这为算法与硬件的协同优化提供了现实可行的计算基础。

上述四个技术方向与ECC设计的三大困境构成精准映射：GNN解决码结构学习的专家依赖，神经算子突破仿真瓶颈，LLM打通需求与标准之间的语义鸿沟，硬件代理模型弥合算法与硬件之间的割裂。\textbf{PrimeCoDec}正是将这四大技术有机融合，构建覆盖ECC设计全链路的统一自动化框架的系统性尝试。